\section{Polygon binding: N=3, 4, 5 bodies}
\label{sec:polygon_geometry}

\subsection{Setup}

We compare regular $N$-polygon clusters — equilateral triangle ($N=3$),
square ($N=4$), and regular pentagon ($N=5$) — each parameterised by the
nearest-neighbour side length $d$.  For each geometry the effective energy is

\begin{equation}
  E_{\rm eff}(d;\,C,m_{\rm sp}) = \underbrace{\frac{N}{8d^2}}_{\rm ZP}
  + V_2(d;\,C,m_{\rm sp})
  + V_3(d;\,C,m_{\rm sp}),
\end{equation}

with the same zero-point convention ${\rm ZP}=N/(8d^2)$,
two-body sum $V_2 = -C^2\sum_{\rm pairs} e^{-m_{\rm sp}r_{ij}}/r_{ij}$,
and three-body sum $V_3 = -C^3\sum_{\rm triples} I_Y(r_{ij},r_{ik},r_{jk};\,m_{\rm sp})$.

\paragraph{Pentagon geometry.}
For $N=5$ with side $d$ the golden ratio $\varphi = 2\cos(\pi/5)\approx1.618$
gives the two distinct inter-vertex distances:
\begin{align}
  &\text{5 adjacent pairs:}&&r = d, \\
  &\text{5 diagonal pairs:} &&r = \varphi d.
\end{align}
The $\binom{5}{3}=10$ triples split into two classes:
\begin{align}
  &\text{Type A (consecutive, 5):}&&(d,\;d,\;\varphi d), \\
  &\text{Type B (skip-one, 5):}   &&(d,\;\varphi d,\;\varphi d).
\end{align}

\subsection{Binding energies at \texorpdfstring{$C=0.155$}{C=0.155}}

Table~\ref{tab:polygon_binding} compares the classical energy minimum
for $m_{\rm sp}=0.1\,l_{\rm Pl}^{-1}$ at the midpoint of the equilateral
3B-only window.

\begin{table}[h]
\centering
\caption{Binding energies for regular $N$-gons.
$m_{\rm sp}=0.1\,l_{\rm Pl}^{-1}$, $C=0.155$.
$E_{\rm min}^{(2B)}$ uses ZP+$V_2$ only; $E_{\rm min}^{(2B+3B)}$ adds $V_3$.}
\label{tab:polygon_binding}
\begin{tabular}{rllrrrrl}
\hline
$N$ & Geometry & Triples & $E_{\rm min}^{(2B)}$ & $E_{\rm min}^{(3B)}$
    & $d_{\rm eq}$ & $V_3/V_2$ & Bound? \\
    &          & $\binom{N}{3}$ & $[E_{\rm Pl}]$ & $[E_{\rm Pl}]$
    & $[l_{\rm Pl}]$ & & \\
\hline
3 & Equilateral & 1  & $+3.6\times10^{-4}$ & $-7.1\times10^{-3}$ & 4.52 & 1.51 & YES \\
4 & Square      & 4  & $+4.5\times10^{-4}$ & $-6.9\times10^{-2}$ & 2.42 & 2.74 & YES \\
5 & Pentagon    & 10 & $+3.6\times10^{-4}$ & $-2.3\times10^{-1}$ & 1.69 & 3.76 & YES \\
\hline
\end{tabular}
\end{table}

All three geometries are in the 3B-only window at $C=0.155$: the
two-body energy is positive (pairs unbound) while the full
two-body+three-body energy is negative (cluster bound).  The binding
energy grows by a factor $\sim32$ from $N=3$ to $N=5$.

\paragraph{Physical origin.}
The dramatic growth comes from two cooperating effects:
\begin{enumerate}
\item \textbf{More triples}: $\binom{N}{3}$ grows as $N^3/6$, giving
  $1\to4\to10$ for $N=3,4,5$.
\item \textbf{Shorter equilibrium separation}: $d_{\rm eq}$ decreases from
  $4.5\to2.4\to1.7\,l_{\rm Pl}$, pushing the cluster deeper into the
  high-$I_Y$ short-distance regime.
\end{enumerate}
The $V_3/V_2$ ratio — a measure of three-body dominance — rises from
$1.5$ to $3.8$, confirming that the cluster becomes progressively more
``three-body driven'' as $N$ increases.

\subsection{Critical coupling thresholds}

\begin{table}[h]
\centering
\caption{Classical critical couplings and 3B-only windows for regular
$N$-gons at $m_{\rm sp}=0.1\,l_{\rm Pl}^{-1}$.
$C_{\rm Pl}=0.282$ lies outside all classical windows for this $m_{\rm sp}$.}
\label{tab:polygon_windows}
\begin{tabular}{rlllll}
\hline
$N$ & Geometry & $C_{\rm cl}^{(2B)}$ & $C_{\rm cl}^{(3B)}$ & $\Delta C$ &
$C_{\rm Pl}$ in? \\
\hline
3 & Equilateral & 0.184 & 0.128 & 0.057 & no \\
4 & Square      & 0.166 & 0.102 & 0.064 & no \\
5 & Pentagon    & 0.159 & 0.091 & 0.068 & no \\
\hline
\end{tabular}
\end{table}

\noindent
\textbf{Key trend:} both thresholds \emph{decrease} with $N$, because
a larger cluster has more 3-body interactions.  Specifically:
\begin{itemize}
\item $C_{\rm cl}^{(3B)}$ decreases from $0.128\to0.102\to0.091$:
  it takes less coupling per triple to bind a larger polygon.
\item $C_{\rm cl}^{(2B)}$ also decreases ($0.184\to0.166\to0.159$)
  because the additional Yukawa pairs in larger polygons aid two-body
  binding.
\item The window width $\Delta C$ \emph{increases} slightly:
  $0.057\to0.064\to0.068$.
\end{itemize}

The 3B-only window \emph{shifts downward} as $N$ increases.
At $m_{\rm sp}=0.1\,l_{\rm Pl}^{-1}$, the window lies in
$C\in(0.091,\,0.184)$ for any of the three polygon types,
but $C_{\rm Pl}=0.282$ lies well above $C_{\rm cl}^{(2B)}$ for all $N$.
This means that at the Planck coupling \emph{all} polygon geometries
have their pairs classically bound; the purely-three-body classical window
does not extend to $C_{\rm Pl}$ for $m_{\rm sp}=0.1$.

\subsection{Pentagon Feynman integrals}

At the pentagon equilibrium ($d_{\rm eq}=1.69\,l_{\rm Pl}$, $C=0.155$):
\begin{align}
  I_Y^{(A)}(d,d,\varphi d) &= 10.16, \\
  I_Y^{(B)}(d,\varphi d,\varphi d) &= 8.77, \\
  \sum_{\rm triples} I_Y &= 5\times10.16 + 5\times8.77 = 94.7.
\end{align}
For comparison, the equilateral triangle at its equilibrium
($d_{\rm eq}=4.52\,l_{\rm Pl}$) gives $I_Y(d,d,d)=4.10$.
The pentagon integral is larger primarily because $d_{\rm eq}$ is smaller
(shorter-range Yukawa is stronger at close range).

\subsection{Comparison at \texorpdfstring{$C=C_{\rm Pl}$}{C=C\_Pl}}

At $C_{\rm Pl}=0.282$, all pair energies are negative; there is no
3B-only classical window.  However, the absolute binding energy grows
dramatically:

\begin{table}[h]
\centering
\caption{Binding energies at $C=C_{\rm Pl}=0.282$,
$m_{\rm sp}=0.1\,l_{\rm Pl}^{-1}$.  All geometries are classically bound
(pairs also bind), so no Efimov-type 3B-only window exists at this coupling.}
\label{tab:polygon_cpl}
\begin{tabular}{rlll}
\hline
$N$ & Geometry & $E_{\rm min}^{(2B)}$ $[E_{\rm Pl}]$ &
$E_{\rm min}^{(2B+3B)}$ $[E_{\rm Pl}]$ \\
\hline
3 & Equilateral & $-0.018$ & $-0.265$ \\
4 & Square      & $-0.051$ & $-1.350$ \\
5 & Pentagon    & $-0.096$ & $-3.729$ \\
\hline
\end{tabular}
\end{table}

The 3-body contribution at $C_{\rm Pl}$ dominates: pentagon binding
($-3.73\,E_{\rm Pl}$) is $14\times$ stronger than the equilateral
($-0.27\,E_{\rm Pl}$), while the 2-body energy ratio is only $5\times$.

\subsection{Summary}

\begin{table}[h]
\centering
\caption{Polygon geometry comparison summary, $m_{\rm sp}=0.1\,l_{\rm Pl}^{-1}$.}
\label{tab:polygon_summary}
\begin{tabular}{llll}
\hline
Property & $N=3$ & $N=4$ & $N=5$ \\
\hline
\# pairs                   & 3    & 6    & 10 \\
\# triples                 & 1    & 4    & 10 \\
$C_{\rm cl}^{(3B)}$        & 0.128 & 0.102 & 0.091 \\
$C_{\rm cl}^{(2B)}$        & 0.184 & 0.166 & 0.159 \\
$\Delta C$ window          & 0.057 & 0.064 & 0.068 \\
$d_{\rm eq}$ ($C=0.155$)   & $4.5\,l_{\rm Pl}$ & $2.4\,l_{\rm Pl}$ & $1.7\,l_{\rm Pl}$ \\
$E_{\rm bind}$ ($C=0.155$) & $7\times10^{-3}$ & $6.9\times10^{-2}$ & $2.3\times10^{-1}$ \\
$V_3/V_2$ ratio            & 1.51 & 2.74 & 3.76 \\
\hline
\end{tabular}
\end{table}

\noindent
\textbf{Main conclusion:} TGP three-body forces create bound polygon
clusters for all $N=3,4,5$ in an Efimov-like 3B-only window.
As $N$ increases: (i) the binding energy grows superlinearly ($\sim N^3$),
(ii) the equilibrium separation shrinks, (iii) the window shifts to lower
$C$, and (iv) three-body forces become progressively more dominant.
The pentagon cluster is $\sim32\times$ more tightly bound than the
equilateral triangle at the same coupling constant.
